\section{Motivation And Workloads}
\label{sec:workload-characterization}

A VFS name-resolution extension targets a narrow workload pattern in which the
source system changes which existing object a pathname should select, but the
operation does not need a new filesystem object model.
The source-system evidence below identifies that pattern and separates it from
the broader behavior that should remain with agent runtimes, environment
builders, services, FUSE filesystems, or custom filesystems.

\subsection{State-Dependent Path Views}

We call a transition a state-dependent path view when a system state change
alters which existing object a pathname denotes, or whether the object is
visible, while ordinary lower-filesystem semantics remain unchanged. This
definition is intentionally narrower than "programmable filesystem." It
excludes synthetic file contents, custom metadata persistence, distributed
directory indexing, data-path mediation, write-conflict resolution, and
cross-path transactional snapshots.

Within this definition, the transition is the workload state change, the policy
is the eBPF decision logic, and the action is the bounded result that the kernel
validates for lookup or directory enumeration.

The distinction matters because the natural implementation boundary changes.
If the system needs synthetic contents or custom metadata persistence, a FUSE
service, custom filesystem, metadata service, or application/runtime mechanism
is the right boundary. If the state change only affects lookup or directory
visibility, a VFS name-resolution policy can be enough. We call the remaining
source-system behavior non-name-resolution responsibilities.

\subsection{Workload-Selection Evidence}

Public systems provide different kinds of evidence: reproduced runs, selected
replays, source inspection, and paper-derived oracles. We use them to choose
representative KVM workloads and to state each oracle's limits, not as a
separate contribution.

A source oracle is the correctness condition extracted from the original source
system. A same-oracle KVM run applies that condition unchanged to \namei and its
comparison rows.

The first source family is agent/workspace state. AgentFS is the primary source
because its public implementation exposes workspace lifecycle traces
from which path-selection and visibility oracles can be extracted~\cite{agentfs_repo}.
BranchFS, Sandlock, and YoloFS provide neighboring branch, sandbox, staging,
and stackable-filesystem context rather than separate evaluation
workloads~\cite{branchfs_repo,sandlock_repo,yolofs_repo}.

The second source family is environment/cache state. SWE-Factory-Gym or
MEnvData-SWE are source systems because their public artifacts provide
Docker-backed build/test oracles~\cite{swe_factory_repo,menvagent_repo}.
SWE-rebench V2 supplies broader benchmark context~\cite{swe_rebench_v2_repo}.
These sources are useful only when the oracle can be reduced to selecting or
hiding existing cache or environment objects while the source evaluator keeps
build/test semantics.

Service/config sources remain conditional. Projected configuration, secret,
certificate, or fixture views motivate the problem, but a service/config
workload belongs in the evaluation only when the source behavior depends on
lookup-time object selection rather than projected-volume materialization or
service reload alone~\cite{kubernetes_projected_volumes}.

Because each evidence source supports claims at different strengths, only
reproduced source behavior can define a final KVM workload oracle. Source
inspection can motivate workload shape, and paper-derived behavior can define a
boundary to test, but the evaluation reports only rows that map to a concrete
same-oracle KVM run.

\begin{table}[t]
\centering
\scriptsize
\begin{tabular}{@{}L{0.2\linewidth}L{0.25\linewidth}L{0.25\linewidth}L{0.2\linewidth}@{}}
\hline
Example & State and operation & \namei action and delegated behavior & Oracle \\
\hline
Agent workspace lifecycle & branch or checkpoint state changes lookup and
readdir under a workspace path & select or hide existing lower objects while
the source runtime keeps COW, checkpoint, sandbox, writes, and lifecycle
orchestration & final tree, lookup/readdir coherence, and preserved
lower-filesystem permissions \\
Environment/cache transition & cache epoch marks a local object hit, miss,
stale, corrupt, or behind an epoch update & select verified local or canonical
fallback, or hide a stale/corrupt local object while cache population, build/test
execution, and content validation stay outside the hook & source build/test
result and stale/corrupt non-use \\
\hline
\end{tabular}
\caption{Representative state-dependent path-view examples. The \namei
evidence is the lookup/readdir action, and non-name-resolution behavior remains
with the source runtime or lower filesystem.}
\label{tab:path-view-examples}
\end{table}

Three design constraints follow from the source systems. First, each workload
carries its own state, such as a branch identifier, checkpoint, cache epoch,
environment epoch, service epoch, or secret epoch. Policy therefore must run at
the affected lookup or directory-enumeration operation and remain scoped per
workload.

Second, the path-view portion of each source system is narrower than the full
system. Many rows also include runtime orchestration, service reload,
synthetic contents, or storage semantics that remain separate responsibilities.
The narrower path-view scope implies that the kernel and lower filesystem must
keep data, permission, and persistence ownership.

Third, directory visibility and lookup selection appear together in the source
oracles. Lookup and readdir therefore need one coherent decision contract.

Together, these source-system observations set the paper's comparison rule:
(1) source behavior supplies the oracle, (2) feature-equivalent FUSE supplies
the direct programmable-policy cost comparison, and (3) custom or stackable
filesystems define the broader safety and ownership boundary. Materialized
namespace mechanisms remain related-work and background comparisons rather than
the central experiment. The same rule drives
the design requirements in Section~\ref{sec:design} and the RQ-organized
evaluation in Section~\ref{sec:evaluation}.

\subsection{Selected Workloads}

This evidence filter leaves two primary representative workloads. One workload
derives from a complete AgentFS-style workspace lifecycle whose
name-resolution decisions select or hide branch-visible objects while COW,
checkpoint, sandbox, and write behavior remain responsibilities of the source
runtime.

The second workload is an environment/cache transition that chooses existing
verified local or canonical objects, and can select regenerated objects only
after the source evaluator has materialized them. It hides stale or corrupt
local objects while the source evaluator and lower filesystem keep cache
population and build/test semantics. This workload exercises file-object cache
choices, which require final-file target selection.

A workload becomes reportable only after a reviewed KVM same-oracle run.
Service/config is included only when a concrete source oracle depends on
lookup-time object selection. These workloads instantiate the boundary studied
by the paper: policy at lookup and directory enumeration, with
non-name-resolution behavior left to the source system or lower filesystem.
